\documentclass[10pt,a4paper]{amsart}
\usepackage[margin=19mm]{geometry}
\usepackage{amsmath,amssymb,mathtools}
\usepackage[scr=rsfs,cal=euler]{mathalfa}
\usepackage{xfrac}
\usepackage[unicode,hidelinks,bookmarks=false]{hyperref}
\newcommand{\ie}{i.e.\ }
\newcommand{\cf}{cf.\ }
\newcommand{\resp}{respectively\ }
\newcommand{\Subsection}{\subsection}
\providecommand{\nobreakdash}{\nobreak\hbox{-}}
\setlength{\emergencystretch}{2em}
\multlinegap0pt
\usepackage{amssymb}
\usepackage{xcolor}
\newtheorem{lemma}{Lemma}
\newcommand*{\CC}{\mathbb C}
\DeclareMathOperator{\End}{End}
\DeclareMathOperator{\U}{Isom}
\newcommand*{\abs}[1]{\lvert#1\rvert}
\DeclareMathOperator{\id}{id}
\newcommand*{\norm}[1]{\lVert#1\rVert}
\newcommand*{\normop}[1]{\lVert#1\rVert_\text{op}}
\newcommand*{\normopU}[1]{\lVert#1\rVert_{U, \text{op}}}
\DeclareMathOperator{\trace}{trace}
\title{A note about Jordan's bound on the size of finite linear groups}
\author{Peter M\"uller}
\date{Source: arXiv:2603.15813v1 (CC BY 4.0)}
\begin{document}
\maketitle

\noindent\emph{Source and licence.} A note about Jordan's bound on the size of finite linear groups. Version arXiv:2603.15813v1 (CC BY 4.0), distributed by arXiv under the Creative Commons Attribution 4.0 International licence, http://creativecommons.org/licenses/by/4.0/, as recorded in the arXiv licence metadata for this exact version and shown on the abstract page https://arxiv.org/abs/2603.15813v1. Full licence notice: \texttt{licenses/arxiv-2603.15813-CC-BY-4.0.txt}.\par\smallskip

\noindent\textit{Editorial scope.} Counted unit: the article's lemma l:xy=yx (source0.tex lines 213--223). The article's complete original proof of it is delivered with it (source0.tex lines 224--294, one \texttt{proof} environment of 71 lines). No mathematics is written, repaired or re-derived: the statement and the proof are the article's own lines, in the article's own notation, and no retained line is substituted or re-flowed.  No further source-local context is needed: every cross-reference of every retained line resolves inside the two delivered spans.  Nothing was omitted from any retained span, so the package carries no omission record.  No retained line contains a citation, so the wrapper carries no bibliography and no bibliography entry.  No retained line contains a figure environment, an image-inclusion command or drawing code, so no image is delivered and the package declares no asset dependency.  Editorial formatting only, in addition: an AMS \texttt{amsart} preamble that restates the article's own declarations and macros used by the retained lines, namely the environments \texttt{lemma} on separate counters, together with the standard packages those lines need.  The retained environments print in this document as Lemma 1 in order of appearance, whereas the article prints the same items with the numbering of its own full document.  The complete native source of the article as distributed in the arXiv e-print for this exact version is bundled unchanged as \texttt{sources/196501/source0.tex}.
\begin{lemma}\label{l:xy=yx}
  Let $V$ be a finite-dimensional complex vector space with norm
  $\norm{\cdot}$ and $\id_V$ be the identity map on $V$. Then the
  following holds:
  \begin{itemize}
  \item[(a)] If $\normop{\id_V-x}<1$ for $x\in\End(V)$, then $\trace
    x\ne0$.
  \item[(b)] Suppose that $x,y\in\U(V)$ generate a finite group. If
    $\normop{\id_V-x}<1/2$ and $\normop{\id_V-y}<1/2$, then $xy=yx$.
  \end{itemize}
\end{lemma}

\begin{proof}
  \begin{itemize}
  \item[(a)] If $\lambda$ is an eigenvalue of $x$, then $1-\lambda$ is
    an eigenvalue of $\id_V-x$, hence
    $\abs{1-\lambda}\le\normop{\id_V-x}<1$.

    Let $\lambda_1,\dots,\lambda_n$ be the eigenvalues of $x$,
    where $n=\dim V\ge1$. Then
    \[
      n %
      > \sum\abs{1-\lambda_i}%
       \ge \abs{\sum(1-\lambda_i)}%
      =\abs{n-\trace x},
    \]
    hence $\trace x\ne0$.
  \item[(b)] We prove the claim by induction on $\dim V$. The base
    case $\dim V=1$ is immediate. Also, clearly $xy=yx$ if $x$ is a
    scalar map. Thus suppose that $x$ is not a scalar map. Let $z$ be
    a non-scalar element of the finite group generated by $x$ and $y$
    for which $\normop{\id_V-z}$ is minimal. From
  \begin{align*}
    \id_V-x^{-1}z^{-1}xz%
    &=x^{-1}z^{-1}(zx-xz)\\
    &=x^{-1}z^{-1}\left((\id_V-z)(\id_V-x)-(\id_V-x)(\id_V-z)\right)
  \end{align*}
  we obtain
  \begin{align*}
    \normop{\id_V-x^{-1}z^{-1}xz}%
    &\le \normop{(\id_V-z)(\id_V-x)-(\id_V-x)(\id_V-z)}\\
    &\le2 \cdot\normop{\id_V-x}\cdot\normop{\id_V-z}\\
    &< \normop{\id_V-z}.
  \end{align*}
  By the minimal choice of $z$ we see that $x^{-1}z^{-1}xz$ is a
  scalar map $\lambda\cdot\id_V$, hence
  \[
    z^{-1}xz=\lambda x
  \]
  for some $\lambda\in\CC$.
  
  Taking the trace gives $\lambda=1$ because $\trace x\ne0$ by (a).

  So $z$ commutes with $x$, and for the same reason $z$ commutes with
  $y$.

  As $z$ has finite order, $V$ is a direct sum of the eigenspaces of
  $z$. These eigenspaces are proper subspaces of $V$, because $z$ is
  not a scalar map. Moreover, as the elements $x$ and $y$ commute with
  $z$, they map each eigenspace of $z$ to itself. Thus in order to
  show that $x$ and $y$ commute, we need only show that the
  restrictions $x_U$ and $y_U$ of $x$ and $y$ to $U$ commute for each
  eigenspace $U$ of $z$.

  Thus the lemma follows by induction on $\dim V$ once we verify that
  $x_U$ and $y_U$ satisfy the assumptions of the lemma.

  Of course, $x_U$ and $y_U$ generate a finite group.

  Let $\norm{\cdot}_U$ be the restriction of $\norm{\cdot}$ to $U$,
  and ${\normopU{\cdot}}$ be the corresponding operator norm on
  $\End(U)$.

  Then $x_U, y_U\in\U(U)$ with respect to
  $\norm{\cdot}_U$. Furthermore, the definition of the operator norm
  shows that $\normopU{\id_U-a_U}\le\normop{\id_V-a}$ for each
  $a\in\End(V)$ which maps $U$ to itself. In particular,
  $\normopU{\id_U-x_U}<1/2$ and $\normopU{\id_U-y_U}<1/2$.

  Thus all the assumptions of the lemma are satisfied for $x_U$ and
  $y_U$, so $x_U$ and $y_U$ commute by the induction hypothesis.
\end{itemize}
\end{proof}

\end{document}
